% -*- TeX-master: "../main.tex" -*-
\chapter{Dostępność}
Aby z naszego narzędzia mogła korzystać szeroka gama osób, wprowadziliśmy ustawienia dostępności, służące do dostosowania w zależności od potrzeb aplikacji. Zdajemy sobie sprawę z wielu wad genetycznych i różnych nieprzyjemności ograniczających możliwości użytkowników. Poniżej sporządziliśmy listę opcji, które sporządziliśmy dla naszego programu.

\section{Daltonizm}

Dużą częścią aplikacji jest kod, którego różne komponenty są kolorowane w zależności od ich zadania. Wiąże się to z możliwością problemu w odczycie dla użytkowników z wadami wzrokowymi, konkretniej z wadami uniemożliwiającymi odróżnianie niektórych barw. Na potrzeby tej grupy stworzyliśmy specjalny motyw na podstawie palety Wong (schemat \ref{tab:dalt}), która celuje w umożliwienie rodzajom daltonizmu takim protanopia, deuteranopia, tritanopia czy monochromatyzm czytelny odczyt.

\section{Tryb jednoręki}
Inną grupą osób z ograniczeniami, dla których chcieliśmy dostosować aplikację, są osoby, które nie są w stanie korzystać z klawiatury oburącz. Na ich potrzeby stworzyliśmy tryb jednoręki z wyborem na lewą lub prawą rękę. Po włączeniu klawiatura użytkownika przesuwana jest w wybraną stronę oraz lekko ścieśniana, co ułatwia wprowadzanie znaków przy pomocy wyłącznie jednej ręki.
